\documentclass[11pt]{article}

% ---------------------------------------------------------------------------
% Packages (all standard; compiles with a normal pdflatex install)
% ---------------------------------------------------------------------------
\usepackage[utf8]{inputenc}
\usepackage[T1]{fontenc}
\usepackage[margin=1in]{geometry}
\usepackage{amsmath,amssymb}
\usepackage{booktabs}
\usepackage{xcolor}
\usepackage{listings}
\usepackage{enumitem}
\usepackage{caption}
\usepackage[hidelinks]{hyperref}

% ---------------------------------------------------------------------------
% Code listing style
% ---------------------------------------------------------------------------
\definecolor{codegray}{rgb}{0.95,0.95,0.95}
\definecolor{codekw}{rgb}{0.10,0.10,0.60}
\definecolor{codecmt}{rgb}{0.30,0.55,0.30}
\definecolor{codestr}{rgb}{0.65,0.20,0.20}
\lstdefinestyle{py}{
  language=Python,
  backgroundcolor=\color{codegray},
  basicstyle=\ttfamily\footnotesize,
  keywordstyle=\color{codekw}\bfseries,
  commentstyle=\color{codecmt}\itshape,
  stringstyle=\color{codestr},
  numbers=left, numberstyle=\tiny\color{gray}, numbersep=6pt,
  breaklines=true, frame=single, framesep=4pt, showstringspaces=false,
  captionpos=b, xleftmargin=2em,
}
\lstset{style=py}

\title{\textbf{Benchmarking Multi-Token-Prediction Models\\
Across Three Decoding Regimes}\\[4pt]
\large Standard, Verified Speculative, and Non-Verification Parallel Decoding\\
on GSM8K, HumanEval, MMLU, and ARC-Challenge}

\author{\textit{Author:} \texttt{<your name>} \\ \textit{Internship Project (TI)}}
\date{27 May 2026}

\begin{document}
\maketitle

\begin{abstract}
This report documents the design and implementation of a benchmark that
measures how the \emph{verification layer} of multi-token-prediction (MTP)
decoding trades off generation \emph{throughput} against task \emph{accuracy}.
We evaluate mid-size language models under three execution modes---standard
single-token decoding, verified (lossless) speculative decoding, and
non-verification parallel decoding---across four reasoning benchmarks: GSM8K
(math), HumanEval (code), MMLU (knowledge), and ARC-Challenge (science). We
explain why the two models named in the original brief
(\texttt{facebook/multi-token-prediction} and
\texttt{unsloth/Qwen3.5-9B-MTP-GGUF}) were substituted with openly loadable
transformers-format models, and we give a detailed technical justification for
emulating the MTP heads with a small \emph{draft model} rather than driving the
target model's native NEXTN heads. The accompanying code is intentionally
simple---plain greedy decode loops behind a single \texttt{propose($\cdot$)}
abstraction---so that the only thing that changes between the three regimes is
how proposed tokens are accepted.
\end{abstract}

\tableofcontents
\bigskip

% ===========================================================================
\section{Introduction}
% ===========================================================================
Large language models normally generate text one token per forward pass. This
is wasteful: a single forward pass already computes a rich representation, yet
it commits to only one output token. \emph{Multi-token prediction} (MTP) adds
auxiliary heads that predict several future tokens from the same pass, and
\emph{speculative decoding} turns those guesses into speedups by drafting
several tokens cheaply and then verifying them with the strong model in
parallel.

The central question of this study is: \textbf{what does the verification step
buy us, and what does removing it cost?} Verification guarantees that the output
is identical to ordinary greedy decoding (it is \emph{lossless}), but it caps
the achievable speedup. If we drop verification and simply accept the cheap
multi-token guesses, generation is faster still---but errors compound and the
generated reasoning can lose coherence. We quantify exactly this trade-off.

\paragraph{Objective.} For each (model $\times$ decoding mode $\times$ dataset)
combination, measure two quantities:
\begin{itemize}[nosep]
  \item \textbf{Accuracy} --- task-specific correctness, to locate the point at
        which unverified text loses logical coherence.
  \item \textbf{Throughput} --- generated tokens per second, to quantify the
        speedup gained by optimizing or removing the verification layer.
\end{itemize}

\paragraph{Contributions of this work.}
\begin{enumerate}[nosep]
  \item A small, readable benchmark harness that runs the same models under
        three decoding regimes and four datasets.
  \item A documented account of the practical gap between the \emph{specified}
        models and what is actually loadable, and the substitutions made.
  \item A reasoned methodological decision to emulate MTP heads with a draft
        model, with the supporting evidence and risk analysis.
\end{enumerate}

% ===========================================================================
\section{Background}
% ===========================================================================
\subsection{Multi-token prediction (MTP)}
A standard causal LM with parameters $\theta$ models
$p_\theta(x_{t+1}\mid x_{\le t})$ and emits one token per step. An MTP model adds
$K$ extra heads so that a single pass over context $x_{\le t}$ produces
predictions for $x_{t+1},x_{t+2},\dots,x_{t+K}$. Head $0$ is the ordinary
next-token head (strong); heads $1\dots K{-}1$ are cheaper look-ahead heads.
This is the design of the Meta MTP models~\cite{gloeckle2024} and the NEXTN
modules in DeepSeek-V3~\cite{deepseekv3} and recent Qwen releases.

\subsection{Speculative decoding and verification}
Speculative decoding~\cite{leviathan2023,chen2023} accelerates generation
without changing its output distribution. A cheap \emph{drafter} proposes a
block of $K$ tokens; the expensive \emph{target} model then scores all $K$
positions in \emph{one} parallel forward pass. Under greedy decoding the
acceptance rule is simply: keep the longest prefix of drafted tokens that equals
the target's own greedy token at each position, then append one free
``bonus'' correction token. Formally, with drafted tokens $d_1,\dots,d_K$ and
target greedy tokens $g_1,\dots,g_K$ at the same positions, the number accepted
is
\begin{equation}
a \;=\; \max\Big\{\, j : d_i = g_i \ \text{for all } i \le j \,\Big\},
\end{equation}
and we emit $d_1,\dots,d_a$ followed by the bonus token $g_{a+1}$. Because every
emitted token agrees with the target's greedy choice, the result is
\textbf{token-for-token identical} to standard greedy decoding---verification is
lossless. The speedup comes from emitting up to $a{+}1$ tokens per target pass
instead of one.

\subsection{The three execution modes under study}
\label{sec:modes}
\begin{description}[style=nextline,leftmargin=1.5em]
  \item[Standard.] Plain greedy decoding: one target forward pass per token.
        The accuracy and (slow) throughput baseline.
  \item[Verified speculative.] Propose $K$, verify in one target pass, accept
        the matching prefix plus bonus token. \emph{Same accuracy as standard}
        (lossless); higher throughput.
  \item[Non-verification parallel.] Propose $K$ and accept \emph{all} of them
        with no verification. Highest throughput, but the target never corrects
        the cheap guesses, so errors accumulate. This is the mode whose accuracy
        we expect to degrade---most on multi-step reasoning where early mistakes
        propagate.
\end{description}

% ===========================================================================
\section{Models: Specification versus Reality}
% ===========================================================================
\subsection{The originally specified models}
The brief named two models:
\begin{enumerate}[nosep]
  \item \texttt{facebook/multi-token-prediction} (7B), the research artifact for
        the Meta MTP paper~\cite{gloeckle2024}.
  \item \texttt{unsloth/Qwen3.5-9B-MTP-GGUF} (9B), a quantized MTP release.
\end{enumerate}

\subsection{Practical obstacles}
Inspecting the published artifacts revealed that neither is directly usable in a
simple \texttt{transformers} script:

\begin{table}[h]
\centering
\caption{Why the specified models are not drop-in loadable.}
\begin{tabular}{@{}p{3.6cm}p{4.0cm}p{4.6cm}@{}}
\toprule
\textbf{Model} & \textbf{Format / runtime} & \textbf{Obstacle} \\
\midrule
\texttt{facebook/multi-token-prediction} &
PyTorch state-dicts in Llama-2 format; custom \texttt{example\_completion.py}
launched via \texttt{torchrun} &
\emph{Gated} (must accept Meta's research licence); not loadable through
\texttt{AutoModelForCausalLM}; needs Meta's bespoke code. \\
\addlinespace
\texttt{unsloth/Qwen3.5-9B-MTP-GGUF} &
GGUF quantization for \texttt{llama.cpp} &
Not a \texttt{transformers} model; MTP runs only via \texttt{llama.cpp}'s
\texttt{-{}-spec-type draft-mtp}. A different runtime entirely. \\
\bottomrule
\end{tabular}
\end{table}

Crucially, the two live in \emph{incompatible runtimes} (a custom Meta codebase
versus \texttt{llama.cpp}), so a single uniform benchmark cannot load both as
written.

\subsection{Substitution decision}
To keep the study runnable, license-free, and inside one runtime, we substitute
two openly loadable, transformers-format models from the same family
(\texttt{Qwen/Qwen3.5-9B} and \texttt{Qwen/Qwen3.5-4B}). Using one family is
deliberate: a \emph{shared tokenizer} lets a single small draft model
(\texttt{Qwen/Qwen3.5-0.6B}) act as the proposer for both targets. The model
identifiers are isolated in \texttt{config.py} so a genuinely distinct
architecture can be swapped into either slot. (Exact repository names should be
confirmed with \texttt{huggingface-cli}, as release names evolve.)

% ===========================================================================
\section{Why a Draft Model Instead of Native MTP Heads}
\label{sec:why-draft}
% ===========================================================================
This is the most important methodological decision in the study, so we justify
it in detail.

\subsection{The native heads exist}
Inspecting \texttt{Qwen/Qwen3.5-4B} confirms a real MTP module is present. Its
\texttt{config.json} declares \texttt{mtp\_num\_hidden\_layers: 1}, and the
checkpoint contains a complete NEXTN-style head:
\begin{lstlisting}[language=bash,basicstyle=\ttfamily\footnotesize]
mtp.layers.0.self_attn.{q,k,v,o}_proj.weight   mtp.layers.0.mlp.*
mtp.fc.weight        mtp.norm.weight
mtp.pre_fc_norm_hidden.weight   mtp.pre_fc_norm_embedding.weight
\end{lstlisting}
This is the DeepSeek-V3 NEXTN recipe~\cite{deepseekv3}: fuse the trunk's hidden
state $h$ with the embedding of the just-predicted token $e$, project, run one
extra transformer layer, and reuse the tied LM head:
\begin{equation}
\label{eq:nextn}
x \,=\, W_{\text{fc}}\,\big[\,\mathrm{RMSNorm}_h(h)\,;\,\mathrm{RMSNorm}_e(e)\,\big],
\quad
h' \,=\, \mathrm{Layer}(x),
\quad
\hat{t} \,=\, \arg\max\;\mathrm{LMHead}\big(\mathrm{RMSNorm}(h')\big).
\end{equation}
So in principle the native heads \emph{are} the drafter and a separate model is
unnecessary.

\subsection{Why we nonetheless use a draft model}
Four reasons, in order of importance:

\paragraph{(1) \texttt{transformers} does not run the heads by default.}
The model class (\texttt{Qwen3\_5ForConditionalGeneration}) \emph{loads} the
\texttt{mtp.*} weights but does not invoke them during \texttt{generate()}. The
NEXTN module is driven only by serving frameworks---vLLM
(\texttt{qwen3\_next\_mtp}), SGLang (\texttt{NEXTN}), and \texttt{llama.cpp}
(\texttt{draft-mtp}). Reaching the heads from a plain Python loop requires
hand-driving the module.

\paragraph{(2) Hand-wiring NEXTN is high-risk and unverifiable offline.}
Equation~\eqref{eq:nextn} hides several details that must be exactly right: the
concatenation order inside $W_{\text{fc}}$, the hidden/token alignment shift,
the rotary-position wiring of the extra attention layer, and the carry of hidden
state across multi-step roll-out. The target is also a \emph{hybrid linear-attention
vision-language model}, which complicates the layer's forward signature. A
subtle error here would be silent: in verified mode it merely lowers the
acceptance rate (slower, still correct), but in \textbf{non-verification
parallel mode it directly corrupts the generated tokens---and accuracy is the
headline metric}. We could not validate the exact wiring without running it on
the target hardware.

\paragraph{(3) Faithful proxy for the quantity under study.}
The study measures the \emph{verification trade-off}, not a specific head
implementation. What matters is having (i) a strong next-token source and (ii) a
cheaper multi-token look-ahead source whose unverified errors we can observe. A
small draft model provides exactly this contrast and is a standard, well-understood
speculative-decoding drafter~\cite{leviathan2023,chen2023}. The three decode
loops are written against one interface, so the conclusions about
verification are independent of whether the look-ahead comes from a draft model
or native heads.

\paragraph{(4) Simplicity and reproducibility.}
The brief asked for simple code that is easy to understand and certain to run.
The draft-model proposer is a handful of lines, has no model-specific assumptions,
and works for any pair of tokenizer-compatible models. It cannot silently
produce wrong physics.

\subsection{Design that keeps the door open}
The decode loops depend only on a \texttt{propose(ids, k)} call. Replacing the
draft model with a validated native-NEXTN proposer is a drop-in change that
leaves the three modes untouched. For trustworthy \emph{native-head} numbers on
the two verifiable modes (standard and verified speculative), vLLM's
\texttt{qwen3\_next\_mtp} is the maintained reference; it cannot, however,
produce the non-verification parallel mode, which is precisely why the
draft-based custom loop remains necessary for the full study.

% ===========================================================================
\section{Methodology}
% ===========================================================================
\subsection{One abstraction, three modes}
All three regimes share a single proposer. The draft proposer returns a block in
which the first token is the strong target prediction and the remainder come
from the draft model:
\begin{lstlisting}[caption={Draft proposer: strong first token, cheap tail.}]
def propose(self, ids, k):
    guesses = []
    cur = ids
    for _ in range(k):
        logits = self.model(cur).logits[:, -1, :]   # draft model
        tok = int(logits.argmax(-1))
        guesses.append(tok)
        cur = torch.cat([cur, tok], dim=1)
    return guesses
\end{lstlisting}

The modes differ only in how the block is consumed:
\begin{lstlisting}[caption={Verified speculative acceptance (lossless).}]
draft_toks = proposer.propose(ids, k)             # cheap block
logits = target(ids + draft_toks).logits          # ONE target pass
target_preds = logits[anchor : anchor+k+1].argmax(-1)
accepted = longest_matching_prefix(draft_toks, target_preds[:k])
emit(draft_toks[:accepted] + [target_preds[accepted]])  # + free bonus token
\end{lstlisting}
\begin{lstlisting}[caption={Non-verification parallel (lossy).}]
anchor_tok = target_next_token(ids)   # one strong token, head 0
emit(anchor_tok)
guesses = proposer.propose(ids, k-1)  # cheap look-ahead
emit(guesses)                         # accepted blindly, NO verification
\end{lstlisting}
The loops use full forward passes (no KV cache) for readability; see
\S\ref{sec:limitations}.

\subsection{Datasets and tasks}
All four datasets are evaluated \emph{generatively} (the model writes an answer
in text) so that every decode mode is exercised identically.

\begin{table}[h]
\centering
\caption{Benchmark datasets and scoring.}
\begin{tabular}{@{}llll@{}}
\toprule
\textbf{Dataset} & \textbf{Domain} & \textbf{Source / split} & \textbf{Scoring} \\
\midrule
GSM8K & grade-school math & \texttt{gsm8k/main}, test & exact match on final integer \\
HumanEval & Python coding & \texttt{openai\_humaneval}, test & unit-test pass@1 (executed) \\
MMLU & general knowledge & \texttt{cais/mmlu/all}, test & extracted choice letter \\
ARC & science reasoning & \texttt{ai2\_arc/ARC-Challenge}, test & extracted choice letter \\
\bottomrule
\end{tabular}
\end{table}

\paragraph{Prompting.} Each example is wrapped with the model's chat template
(when present). Math and multiple-choice prompts request a parseable final line
(``\texttt{Answer: <value>}''); the code prompt requests a fenced Python block.

\paragraph{Answer extraction.} GSM8K: last integer (after \texttt{\#\#\#\#} in the
gold). MMLU/ARC: an explicit ``Answer: X'' letter, else the last standalone
A--D. HumanEval: the fenced code block, completed with the signature if needed,
then executed against the official tests.

\subsection{Metrics}
\textbf{Accuracy} is the mean task score over the evaluated examples.
\textbf{Throughput} is total generated tokens divided by total wall-clock
generation time on the target model:
\begin{equation}
\text{throughput} \;=\; \frac{\sum_i (\text{new tokens})_i}{\sum_i (\text{seconds})_i}
\quad [\text{tokens/s}].
\end{equation}
We also record the number of expensive target forward passes, which is the
mode-independent driver of the speedup.

\subsection{Experimental matrix}
The full sweep is $2\ \text{models} \times 3\ \text{modes} \times 4\
\text{datasets} = 24$ cells, each reporting accuracy and throughput. A
\texttt{-{}-limit} flag caps examples per dataset for fast iteration.

% ===========================================================================
\section{Implementation}
% ===========================================================================
\subsection{Code structure}
\begin{table}[h]
\centering
\caption{Repository layout.}
\begin{tabular}{@{}ll@{}}
\toprule
\textbf{File} & \textbf{Responsibility} \\
\midrule
\texttt{environment.yml} & Conda environment (Python 3.11, CUDA-13.0 PyTorch, HF stack) \\
\texttt{config.py}       & All knobs: models, datasets, modes, limits, $K$ \\
\texttt{mtp\_decoding.py} & Model loading, draft proposer, the three decode loops \\
\texttt{datasets\_tasks.py} & Prompt building, answer extraction, scoring \\
\texttt{benchmark.py}    & Driver: sweeps the matrix, prints and saves results \\
\bottomrule
\end{tabular}
\end{table}

\subsection{Environment and hardware}
Experiments target a single large-memory NVIDIA GPU (128~GB), so both targets
load in \texttt{bfloat16} without quantization and speculative decoding fits
comfortably. The environment pins PyTorch built for \textbf{CUDA 13.0}
(\texttt{cu130}, requiring \texttt{torch>=2.10}). Because the \texttt{cu130}
wheel index omits a \texttt{cuda-bindings} dependency, the install uses
\texttt{-{}-extra-index-url} rather than \texttt{-{}-index-url}. CUDA drivers are
backward compatible, so a 13.0 driver also runs older \texttt{cu128}/\texttt{cu126}
wheels if needed.

\subsection{Running the benchmark}
\begin{lstlisting}[language=bash,basicstyle=\ttfamily\footnotesize]
conda env create -f environment.yml
conda activate mtp-bench
huggingface-cli login                 # if any chosen repo is gated

python benchmark.py --limit 10        # quick smoke test
python benchmark.py --limit 100 --allow-code-exec   # full, with code execution
python benchmark.py --models Qwen3.5-9B --tasks gsm8k --modes standard parallel
\end{lstlisting}
Results print as a table and are saved to \texttt{results/results.json}.

% ===========================================================================
\section{Expected Outcomes and Hypotheses}
% ===========================================================================
\begin{description}[style=nextline,leftmargin=1.5em]
  \item[H1 (lossless verification).] \texttt{standard} and \texttt{speculative}
        yield \emph{identical} accuracy on every dataset; \texttt{speculative}
        is faster.
  \item[H2 (throughput ordering).] Throughput increases
        \texttt{standard} $<$ \texttt{speculative} $<$ \texttt{parallel}, since
        each step removes target work.
  \item[H3 (coherence collapse).] \texttt{parallel} accuracy falls below the
        other two, with the largest drop on multi-step reasoning (GSM8K, ARC),
        because unverified errors compound; single-shot extraction tasks degrade
        more gently.
\end{description}
The study's deliverable is the curve of accuracy versus throughput that these
three points trace out for each model and dataset.

% ===========================================================================
\section{Limitations and Threats to Validity}
\label{sec:limitations}
% ===========================================================================
\begin{itemize}[nosep]
  \item \textbf{Substituted models.} Results characterize the substituted Qwen3.5
        models, not the originally specified Meta/GGUF checkpoints.
  \item \textbf{Emulated heads.} Look-ahead comes from a draft model, not the
        target's native NEXTN heads (\S\ref{sec:why-draft}). Absolute speedups
        differ from a fused MTP implementation; the \emph{direction} and the
        accuracy trade-off are faithful.
  \item \textbf{No KV cache.} Decode loops use full forward passes for clarity,
        so absolute tokens/s is lower than production; \emph{relative} mode
        comparisons remain valid.
  \item \textbf{Small default \texttt{-{}-limit}.} Statistically meaningful
        accuracy requires raising the example count.
  \item \textbf{Code execution.} \texttt{-{}-allow-code-exec} runs
        model-generated Python (subprocess + timeout). This is not a security
        sandbox; run it only inside a container or VM.
\end{itemize}

% ===========================================================================
\section{Future Work}
% ===========================================================================
\begin{itemize}[nosep]
  \item Add a validated native-NEXTN proposer (self-checked by acceptance rate)
        as a drop-in replacement for the draft model.
  \item Cross-check the two verifiable modes against vLLM's
        \texttt{qwen3\_next\_mtp} for trustworthy native-head throughput.
  \item Add KV-cache decoding for realistic absolute throughput.
  \item Sweep the block size $K$ to map the speedup/accuracy frontier of the
        parallel mode.
  \item Include a genuinely distinct second architecture (e.g.\ a DeepSeek/GLM
        NEXTN model) to compare MTP designs, per the original brief.
\end{itemize}

% ===========================================================================
\begin{thebibliography}{9}
% ===========================================================================
\bibitem{gloeckle2024} F.~Gloeckle, B.~Youbi Idrissi, B.~Rozi\`ere, D.~Lopez-Paz, G.~Synnaeve.
``Better \& Faster Large Language Models via Multi-token Prediction.'' arXiv:2404.19737, 2024.

\bibitem{leviathan2023} Y.~Leviathan, M.~Kalman, Y.~Matias.
``Fast Inference from Transformers via Speculative Decoding.'' ICML, 2023.

\bibitem{chen2023} C.~Chen et al.
``Accelerating Large Language Model Decoding with Speculative Sampling.'' arXiv:2302.01318, 2023.

\bibitem{deepseekv3} DeepSeek-AI.
``DeepSeek-V3 Technical Report.'' arXiv:2412.19437, 2024. (Multi-token prediction / NEXTN module.)

\bibitem{cai2024medusa} T.~Cai et al.
``Medusa: Simple LLM Inference Acceleration Framework with Multiple Decoding Heads.'' arXiv:2401.10774, 2024.

\bibitem{cobbe2021} K.~Cobbe et al.
``Training Verifiers to Solve Math Word Problems'' (GSM8K). arXiv:2110.14168, 2021.

\bibitem{chen2021codex} M.~Chen et al.
``Evaluating Large Language Models Trained on Code'' (HumanEval). arXiv:2107.03374, 2021.

\bibitem{hendrycks2021} D.~Hendrycks et al.
``Measuring Massive Multitask Language Understanding'' (MMLU). ICLR, 2021.

\bibitem{clark2018arc} P.~Clark et al.
``Think you have Solved Question Answering? Try ARC.'' arXiv:1803.05457, 2018.
\end{thebibliography}

\end{document}
